\subsection{正top非零の曲がったClifford族とconic判定}
\label{subsec:h3-order2-22-clifford-curved}
以下は全正top非零領域の被覆ではなく、明示的な一族の認識である。
$U=kx+ky$、$V=ku+kv$ とし、任意の二元二次式 $h\in\operatorname{Sym}^2U$ に対して
\[
 B_h=k\langle x,y,u,v\rangle/
 ([x,y],\{x,u\},\{y,u\},\{x,v\},\{y,v\},\{u,v\}-h)
\]
と置く。$h=0$ は既知の非CY符号包絡である。

\begin{proposition}[包絡と正常三次の全核]
\label{prop:h3-order2-22-clifford-normal}
全 $h$ で $B_h$ はnoetherian quadratic Koszul AS4 domainであり、
$\mu_{B_h}|_U=1,\mu_{B_h}|_V=-1$。
$\mu(G)=G$、$Gz=\mu^{-1}(z)G$ を満たす三次式はちょうど
\[
 G=f_3(x,y)+\ell_2(x,y)u^2+\ell_3(x,y)v^2
 \quad(f_3\in\operatorname{Sym}^3U,\ \ell_i\in U)
\]
の八次元空間である。
\end{proposition}
\begin{proof}
$k[x,y][u;-1]$ に、$\sigma(x,y,u)=(-x,-y,-u)$、
$\delta(x)=\delta(y)=0,\delta(u)=h$ のOre変数 $v$ を加える。
$h$ が正側の偶次数なので $\delta$ は関係を保つ。
従って反復Ore拡大のdomain/noetherian/AS正則性が適用できる。
順序 $x<y<u<v$ の四つの三次overlapは全係数で解消し、
PBW基底 $x^iy^ju^kv^l$ とKoszul性を得る。
Ore中山公式を二回用いると $U$ 上は恒等、$u$ 上は負号である。
$v$ の像は $-v+b$ ($b\in kx+ky+ku$) と書けるが、
自己同型 $u\mapsto tu,v\mapsto t^{-1}v$ との可換性が $b=0$ を強制する。

正固有三次は表示した八語と $xuv,yuv$ の計十語。
前八語は直接全normalizer式を満たす。
後二語の $u$ との反交換子の $xu^2v,yu^2v$ 係数はともに $-2$ であり、
その $2\times2$ minorは4。他の八列は零なので全核が確定する。
\end{proof}

以下では $\ell_2,\ell_3$ が独立である開集合を扱う。
正側の座標を取り直して
\[
 G=f_3+xu^2+yv^2,\qquad
 f_3=ax^3+bx^2y+cxy^2+dy^3=xf_2+yg_2,
 \quad f_2=ax^2+bxy+cy^2,\quad g_2=dy^2
\]
とする。$G\ne0$ は上のdomainで正則である。

\begin{theorem}[独立係数chartのH3必要十分条件]
\label{thm:h3-order2-22-clifford-conic}
$S=\mathbb F_1$、$\mathcal V=\OO(F)\oplus\OO\oplus\OO$、
$L=\OO(E+2F)$ とし、$\lambda,\mu,s,t$ の束次数をそれぞれ
$E,E+F,F,F$ とする。次の $L$ 値形式を考える：
\[
 q=\begin{pmatrix}\lambda&0&0\\0&A&h_S/2\\0&h_S/2&B\end{pmatrix},\qquad
 A=-\mu t-\lambda f_2(s,t),\quad B=\mu s-\lambda g_2(s,t),\quad
 h_S=\lambda h(s,t).
\]
$R=B_h/(G)$ は常にdomainであり、そのeven/odd sheafは
$\mathcal C_0(q),\mathcal C_1(q)$ である。
さらに次の三条件は同値である：
\begin{enumerate}
\item $R$ はH3を定める。
\item $q$ のconic bundle総空間が滑らかである。
\item 平面四次式
\[
 F_4=4(-YZ-f_2(X,Y))(XZ-g_2(X,Y))-h(X,Y)^2
\]
は被約で通常node以外の特異点を持たず、$o=[0:0:1]$ 以外の全特異点で
$-YZ-f_2=XZ-g_2=h=0$ が成立する。
\end{enumerate}
$o$ は常に通常nodeで、その二枝は例外曲線上の二つの良いrank-one点に分離する。
\end{theorem}
\begin{proof}
$P=x^2,Q=xy,T=y^2,A_0=u^2,B_0=v^2$ は中心である。
$U_0=A_0+f_2(P,Q,T),V_0=B_0+g_2(P,Q,T)$ とすると
$G=xU_0+yV_0$。
従って中心から作用するscroll環は
\[
 C=k[P,Q,T,U_0,V_0]/(Q^2-PT,PU_0+QV_0,QU_0+TV_0).
\]
その標準表示は $P=\lambda s^2,Q=\lambda st,T=\lambda t^2,
U_0=-\mu t,V_0=\mu s$ であり、$\operatorname{Proj}C=S$。
四変数Clifford行列は正側block
$\left(\begin{smallmatrix}P&Q\\Q&T\end{smallmatrix}\right)$ と負側block
$\left(\begin{smallmatrix}A_0&h/2\\h/2&B_0\end{smallmatrix}\right)$ の直和である。
正側radical $(-t,s)$ を除き $x=s\ell,y=t\ell$ と書けば正確に $q$ が得られる。
このradicalの線型関係は $Qx-Py=0,Tx-Qy=0,U_0x+V_0y=0$ で、
全て $R$ で成立する。逆の局所Clifford表示からも元の六二次関係と $G$ が復元される。
従って局所表示は互いに逆であり、even/odd sheafの同定を得る。
これは $C$ の作用が忠実であることも示す：kernelのsheafは零、
一方 $C$ はdomainで有限長の非零斉次idealを持たない。

頂点の問題を分離する。$B_h$ は中心多項式環
$Z=k[x^2,y^2,u^2,v^2]$ 上、四指数の偶奇による自由基底16個を持つ。
$G$ 倍による自由解消と $G$ の正則性から $R$ は $Z$ 上深さ3。
$C$ は $Z$ の像上有限で $Q^2=PT$、両irrelevant idealのradicalは同じだから、
even/odd部分の局所コホモロジーは次数0,1,2で零である。
従って両部分は上の局所自由Clifford sheafの飽和section moduleであり、
特に $C$ 上torsion-freeである。

$q$ は $\lambda\ne0$ で非零、$E=(\lambda=0)$ 上も
$A=-\mu t,B=\mu s$ が同時零にならず、rank0点はない。
判別は $\lambda(AB-h_S^2/4)$。
$E$ 上の後因子は $-\mu^2st$ なので、二地点 $s=0,t=0$ で $E$ と横断する。
例えば $t=0$ で $B$ は単位、単位方向を除いた形式は
$\operatorname{diag}(\lambda,A-h_S^2/(4B))$ で、両成分は正則parameterである。
局所even Clifford代数は完備量子平面、そのquaternionは
$(\lambda,A-h_S^2/(4B))$。
後者の $\lambda=0$ residueはもう一方のparameterの平方類で非自明である。
従ってglobal generic quaternionもdivisionである。
飽和性によるtorsion-free性と正則なodd方向でのskew Laurent表示から、
定理\ref{thm:h3-order2-global-conic-bundle}のdomain論証と同様に $R$ はdomainとなる。
ここでは同定理の $(1,3)$ 被覆を使用せず、上で個別に再構成した表示を用いた。

基礎束は
$\mathcal C_0=\OO\oplus\OO(-E-F)^{\oplus2}\oplus\OO(-E-2F)$。
従って構造対象は例外的である。正規商のAS Gorenstein parameterは1、
中山は恒等であり、有限中心、飽和性、PBW/Hilbert値とgeneric divisionが揃った。
系\ref{cor:h3-order2-conic-recognition}の証明をこの表示に適用すると、
残るH3条件は全局所orderの大域次元2にちょうど一致する。
rank3はAzumaya、rank2は判別の滑らかさ、rank1は
補題\ref{lem:h3-order2-local-conic-criterion}の横断node条件である。
従って(1)と(2)は同値。

吹き下ろし $(X,Y,Z)=(\lambda s,\lambda t,\mu)$ に対して
$F_4(\lambda s,\lambda t,\mu)=4\lambda\det q$。
$o$ での二次部分は $-4XY$ なので通常nodeである。
判別因子は $E$ と四次曲線のstrict transformの和で、類は $3E+4F$。
$E$ のgeneric点は滑らかなrank2判別、二交点は上で確認した良いrank1点である。
$E$ の外では $\lambda$ は単位なのでrank1は正確に $A=B=h_S=0$。
従って局所判定により(2)と(3)も同値となる。
\end{proof}

\begin{example}[非零Clifford curvatureのH3]
$h=xy$、$G=x^3+y^3+xu^2+yv^2$ とすると
\[
 F_4=-4XYZ^2+4Y^3Z-4X^3Z+3X^2Y^2.
\]
$X=1$ chartの $F_4,\partial_YF_4,\partial_ZF_4$ のGroebner基底は $\{1\}$。
$X=0,Y=1$ では $\partial_ZF_4=4$、従って唯一の特異点は $o$。
上の定理により実際にH3を得る。全係数PBW、正常核minor4、
四次恒等式とこのJacobian判定は次の検証器で厳密再生する：
\[
\texttt{verify\_cubic\_order2\_22\_clifford\_curved.py}.
\]
新規性・独立査読は未確認であり、全正top非零対象の被覆・全矢・族降下はこの節では主張しない。
\end{example}

\begin{proposition}[依存係数境界は全て非H3]
\label{prop:h3-order2-22-clifford-dependent}
この包絡族でH3となる切断は必ず $\ell_2,\ell_3$ が独立である。
従って定理\ref{thm:h3-order2-22-clifford-conic}は、$B_h$ 全体の全H3切断を尽くす。
\end{proposition}
\begin{proof}
両係数が零なら $G$ は二元三次式なので分解し、商はdomainでない。
それ以外の従属係数は、正側を取り直して
$G=f_3+x(\alpha u^2+\beta v^2)$ と書ける。
必要なら $u,v$ を交換して $\beta\ne0$ とする。
H3であると仮定し、中心元 $u^2$ を逆転する。
次数零chartで
\[
 X=xu^{-1},\quad Y=yu^{-1},\quad
 W=vu^{-1}+\tfrac12h(X,Y)
\]
と置くと、正確に
\[
 XY=YX,\quad WX=-XW,\quad WY=-YW.
\]
\[
 F=f_3(X,Y)-\alpha X+\beta X\bigl(W^2-h(X,Y)^2/4\bigr)=0
\]
となる。逆表示は $x=Xu,y=Yu,v=(W-h(X,Y)/2)u$、
$u$ の共役は $X,Y,W$ 全てに負号であり、元の関係を復元する。
特に $G=-Fu^3$ で、付加的な関係のないOre chartである。

このchartは、同じ $F$ を持つ可換affine曲面の有限群cocycle twistである。
実際 $\Gamma=(\mathbb Z/2)^2$、$\deg X=\deg Y=e,\deg W=f$ とし、
$\alpha((i,j),(i',j'))=(-1)^{ji'}$ を用いる。
$F$ は全体が次数 $e$ であり、順序を $X,Y,W$ とした各項のtwist係数は1。
定理\ref{thm:noncy-sign-cubic}の有限同変化・誘導忘却の議論はaffine加群圏にも適用でき、
両環の大域次元は等しい。

$\beta c^2=\alpha$ を満たす $c$ を取ると、可換曲面の点 $(0,0,c)$ では
$F=\partial_XF=\partial_YF=\partial_WF=0$。
従って可換局所環は非正則で、その剰余体の射影次元は無限である。
chartも無限大域次元となる。
しかしH3なら中心開集合への制限もhomological dimension高々2であり、
強次数付き局所化の次数零環としてのこのchartも高々2でなければならない。矛盾。
全六二次関係の逆復元、$G=-Fu^3$、特異点の四恒等式は
\[
\texttt{verify\_cubic\_order2\_22\_clifford\_boundary.py}
\]
で全係数について再生できる。
\end{proof}

\begin{theorem}[曲がったClifford族の全幾何学的矢]
\label{thm:h3-order2-22-clifford-arrows}
全H3係数開集合を
$\mathcal U\subset\mathbb A^3_h\times\mathbb P^7_G$ と書く。
この族のpointed exact幾何学的群oidは
\[
 [\mathcal U/\overline{\mathcal H}](k),\qquad
 \mathcal H=\mathrm{GL}(U)\times((\mathbb G_m)^2\rtimes S_2),\quad
 \overline{\mathcal H}=\mathcal H/\mathbb G_{m,\mathrm{degree}}
\]
である。矢は $M\in\mathrm{GL}(U),\gamma,\delta\ne0$ による
$U\mapsto MU'$ と、
$(u,v)\mapsto(\gamma u',\delta v')$ 又は
$(\gamma v',\delta u')$ のうち、
\[
 h(MU')=\gamma\delta h'(U'),\qquad G(MU',\gamma u',\delta v')\in k^*G'
\]
を満たすもの全て（交換時は二つの負変数も交換）で尽くされる。
表示stackの一般次元は5、一般慣性は $\mu_2$。
$h=0$ は既知の一般次元2・一般慣性 $(\mu_2)^2$ の符号部分族である。
全intrinsic stackの成分性や非被約底上の全降下は主張しない。
\end{theorem}
\begin{proof}
内在的反標準環と二次包絡の回復
（命題\ref{prop:intrinsic-ring-isomorphisms}）により、全矢は包絡の次数付き同型から来る。
中山は実際の $++--$ なので、その二つの固有空間を保つ。
二次関係の $V\otimes V$ への射影は $k(uv+vu)$ であるから、
負側の線型写像はこの非退化対称tensorの直線を保つ。
二変数ではその群は対角変換と交換からなり、上の式が全関係保存の必要十分条件となる。
三次関係についてはその直線を保つことが必要十分であり、次数スカラーだけを除く。
これで特殊係数も削除せず全矢が確定する。

係数独立は開条件、conic総空間のrelative singular incidenceは射影的なので、
必要十分H3条件も開条件である。
実例により $\mathcal U$ は非空で、既約な10次元係数空間の開部分。
有効作用群の次元は $4+2-1=5$ だから一般次元は5。
一般慣性を直接確認する。$\ell_2=x,\ell_3=y$、$abcd\ne0$ と規格化する。
負側交換を含まない安定化では $M=\operatorname{diag}(m,n)$。
$f_3(MU)=c_0f_3(U)$ の四係数から $m=n,c_0=m^3$ となり、
切断の残る二係数から $\gamma^2=\delta^2=m^2$。
$h\ne0$ なら $h(MU)=\gamma\delta h(U)$ は $\gamma\delta=m^2$ を強制する。
従って次数スカラーを除けば恒等と $\mu_B$ の二つだけである。
交換を含む安定化は、切断係数だけから既知の条件 $db^3=ac^3$ を必要とするので一般には存在しない。
これらの一般条件と非空H3開集合は既約係数空間で交わるから、一般慣性は正確に $\mu_2$。
$h=0$ の主張は定理\ref{thm:fixed-sign-classification}に一致する。
\end{proof}
